%! Author = murfe
%! Date = 08.09.2026

\chapter{Methodik}
\label{ch:methodik}
Dieses Kapitel beschreibt das methodische Vorgehen zur Beantwortung der in \cref{sec:problemstellung} formulierten
Forschungs- und Teilfragen.
Dazu wird zuerst das Vorgehen bei der Literaturrecherche dargelegt.
Anschließend folgt die Erläuterung zum Prozess der Anforderungserhebung, der Entwicklung der
Systemarchitektur sowie deren prototypische Umsetzung.
Schließlich wird dargestellt, wie im Rahmen der Evaluation untersucht wird, ob das entwickelte System die
definierten Zusicherungen und Anforderungen einhält.

\section{Literaturrecherche}
\label{sec:literaturrecherche}
Die grundlegende Literaturliste für diese Arbeit wurde durch ein mehrstufiges Verfahren generiert.
Zuerst wurden der Problemkontext und die in \cref{sec:problemstellung} formulierten Forschungs- inklusive Teilfragen nach
den zugrundeliegenden Themen und vorkommenden Schlagworten analysiert.
Diese Liste umfasste Begriffe wie asynchrone Verarbeitung, verteilte Systeme, Fehlertoleranz und Teilausfälle.
Da die vorliegende Arbeit stark architektur- und praxisgetrieben ist, fokussierte sich die systematische Suche primär
auf Grundlagenwerke der Softwarearchitektur und verteilter Systeme sowie auf offizielle technische Dokumentationen.

Ein Teil der Literatur, insbesondere zu Transaktionsverwaltung und Nebenläufigkeitskontrolle in
relationalen Datenbanksystemen (u.\,a. \textcite{saake_datenbanken_2018}), war bereits aus vorangegangenen
Lehrveranstaltungen des Studiums bekannt und wurde auf Relevanz für die Verwendung in dieser Arbeit lediglich erneut geprüft,
statt recherchiert zu werden.

Zur Identifikation von Fachliteratur zu theoretischen Grundlagen wurden zunächst Google Scholar und O'Reilly Media herangezogen.
Für spezifische Problemstellungen, wie etwa typische Ausfallursachen in Cloud-Umgebungen, wurde ergänzend in den
Datenbanken der ACM Digital Library und IEEE Xplore recherchiert.
Der zentrale Suchstring wurde anhand der zuvor erarbeiteten und ins Englische übersetzten Liste an Fachbegriffen abgeleitet
und je nach Plattform syntaktisch angepasst.
Nach Sichtung der ersten Treffer wurde der Suchstring mehrfach angepasst, bis sich schließlich folgender ergab:

(\enquote{partial failure} OR \enquote{unreliable} OR \enquote{network faults}) AND (\enquote{distributed systems}
OR \enquote{cloud}) AND (\enquote{architecture} OR \enquote{design})

Aus der so generierten Ergebnismenge wurde händisch gefiltert.
Zuerst wurden anhand des Titels offensichtlich themenfremde Ergebnisse entfernt.
Bei Fachbüchern erfolgte die Auswahl anhand des Inhaltsverzeichnisses und der Relevanz der Kapitel für die Forschungsfragen
dieser Arbeit.
Bei den wissenschaftlichen Publikationen wurden nach kurzer inhaltlicher Durchsicht der Abstracts nicht relevante Werke aussortiert.
Wichtig für den Einbezug der Literatur war thematische Nähe zur Architektur und zu Fehlern in verteilten Systemen, asynchroner Verarbeitung
oder Fehlertoleranz in verteilten Systemen.

Ausgehend von den so identifizierten Werken wurden mithilfe einer rückwärtsgerichteten Schneeballsuche weitere relevante
Publikationen identifiziert.
Ergänzt wurde die Liste außerdem durch technische Dokumentationen der verwendeten und verglichenen Technologien.

\section{Systementwurf}
\label{sec:systementwurf}
Die Systementwicklung folgt einem schrittweisen, an den Anforderungen orientierten Vorgehen.
Ausgehend von der Problemstellung werden in \cref{ch:anforderungen} zunächst Anforderungen an das zu entwickelnde System
erhoben und daraus konkrete Zusicherungen abgeleitet (vgl.~\ref{item:tz1}).
Diese Zusicherungen beschreiben die Eigenschaften, die das System auch unter Fehlerbedingungen einhalten soll.
Auf dieser Grundlage werden in \cref{ch:umsetzung} Architekturentscheidungen getroffen, die die Erfüllung
der Zusicherungen ermöglichen (vgl.~\ref{item:tz2}).
Die ausgewählten Mechanismen werden anschließend in einem prototypischen System implementiert (vgl.~\ref{item:tz3}).

Damit ergibt sich für den Systementwurf folgende methodische Reihenfolge:
\begin{enumerate}
    \item{Problemstellung}
    \item{Anforderungserhebung}
    \item{Zusicherungen}
    \item{Architekturentwurf}
    \item{prototypische Implementierung}
\end{enumerate}

\subsection{Anforderungserhebung und Ableitung der Zusicherungen}
\label{subsec:anforderungserhebung}
Um die Anforderungen für das System systematisch zu ermitteln, orientiert sich das Vorgehen an den Ermittlungsgegenständen des
Requirements-Engineerings nach \textcite{rupp_requirements-engineering_2021}.
Das Ziel des Systems ist die Entwicklung eines Job-Processing-Systems, das resilient gegenüber Fehlern und Teilausfällen in
verteilten Systemen ist.
Um den Rahmen der Anforderungen dafür einzugrenzen, wird in \cref{sec:systemkontext} der Systemkontext spezifiziert.
Anstelle klassischer Stakeholder-Befragungen dienen in dieser Arbeit die in \cref{sec:abgrenzungFehlermodell}
definierten Fehlerklassen als primäre Anforderungsquelle.
Diese sind eine Teilmenge der in der Literatur beschriebenen typischen Fehlerquellen verteilter Systeme und bilden das für
diese Arbeit betrachtete Fehlermodell ab.

Darauf folgend werden konkrete Zusicherungen abgeleitet.
Diese Verfeinerung hilft, um das zentrale Qualitätskriterium der Prüfbarkeit sicherzustellen.
Eine Zusicherung beschreibt in diesem Sinne eine testbare Eigenschaft, die unabhängig von einem konkreten
Ausführungsablauf beziehungsweise dem Auftreten eines Fehlers zwingend eingehalten werden soll.
Die auf diese Weise formulierten Zusicherungen Z-1 bis Z-9 sind in \cref{sec:zusicherungen} vollständig spezifiziert.
Sie dienen als präzise Korrektheitsbedingungen für das entwickelte System und bilden daher die Grundlage für
den späteren Systementwurf und die Evaluation.

\subsection{Architekturentwurf und prototypische Implementierung}
\label{subsec:architekturentwurfUndImplementierung}
Auf Grundlage der definierten Anforderungen und Zusicherungen werden die Architektur und die zentralen
Verarbeitungsmechanismen des Systems entworfen.
Die Fachliteratur empfiehlt, komplexe Systeme aus mehreren Blickwinkeln zu analysieren und zu dokumentieren \parencite[S.470]{rupp_requirements-engineering_2021},
da die Bildung spezifischer Sichten hilft, das System strukturiert und schrittweise zu entwickeln \parencite[S.469]{rupp_requirements-engineering_2021}.
Das Vorgehen orientiert sich hierfür an dem Vier-Sichten-Modell nach \textcite[S.80 ff.]{starke_effektive_2024}.

Die geforderten Sichten spiegeln sich in der Struktur dieser Arbeit wie folgt wider:
Die Kontextabgrenzung findet bereits im Rahmen der Anforderungserhebung statt und ist entsprechend in \cref{sec:systemkontext} thematisiert.
Die Bausteinsicht wird in \cref{fig:componentDiagram} im Rahmen des Architekturentwurfs dargestellt, während einzelne
Laufzeitsichten in den \cref{sec:architektur} untergeordneten Abschnitten zu finden sind.
Auf eine gesonderte Verteilungssicht wird bewusst verzichtet, da der Prototyp ausschließlich aus einem Client, einem Spring-Boot-Backend
und einer Datenbank besteht.
Eine solche zusätzliche Sicht über den bereits vorhandenen Systemkontext und die Bausteinsicht hinaus würde im Kontext der
Arbeit keinen wesentlichen Mehrwert liefern.

Die resultierende Architektur wird anschließend prototypisch in Java und Spring implementiert.
Da es sich um eine prototypische Umsetzung des Job-Processing-Systems mit Untersuchung der relevanten Fehlertoleranzmechanismen handelt,
wird die eigentliche Jobverarbeitung durch eine simulierte Fachlogik repräsentiert (vgl. \cref{sec:systemkontext}).
Die Implementierung inklusive der relevanten technischen Entscheidungen wird in \cref{sec:implementierung} detailliert beschrieben.

\section{Evaluation}
\label{sec:evaluationsmethodik}
Im Anschluss an die Implementierung untersucht die Evaluation im Rahmen von \ref{item:tz4}, ob die implementierte Lösung
die in \cref{ch:anforderungen} definierten Zusicherungen auch unter gezielt herbeigeführten Fehlern einhält.
Dazu werden die Zusicherungen zunächst in beobachtbare Eigenschaften und konkrete Prüfkriterien überführt.
Auf dieser Grundlage werden geeignete Test- und Nachweismethoden definiert, Fehlersituationen gezielt erzeugt und die
resultierenden Systemzustände ausgewertet.

Die Evaluationsmethodik umfasst damit die folgenden Schritte:
\begin{enumerate}
    \item{Operationalisierung}
    \item{Definition von Tests und Nachweismethoden}
    \item{Fehlerinjektion und Test}
    \item{Auswertung}
\end{enumerate}

\subsection{Operationalisierung der Zusicherungen}
\label{subsec:operationalisierung}
Die in \cref{ch:anforderungen} formulierten Zusicherungen beschreiben zunächst abstrakte Eigenschaften des zu entwickelnden Systems.
Damit diese im Rahmen der Evaluation überprüft werden können, werden sie in konkret beobachtbare Eigenschaften und daraus
ableitbare Prüfkriterien überführt.
Die Operationalisierung legt somit fest, anhand welcher beobachtbaren Systemzustände oder -ereignisse entschieden wird,
ob eine Zusicherung unter einem definierten Fehlerszenario eingehalten wurde.
Hierzu wird für jede Zusicherung zunächst der für den Test erwartete Systemzustand in Form eines Prüfkriteriums formuliert,
das anhand von beobachtbaren Größen wie Zuständen im Job-Store, der Anzahl persistierter Ergebnisse, HTTP-Antworten, ausgeführten
Operationen oder der Anzahl erreichter Verarbeitungsversuche überprüft werden kann.
Die vollständige Zuordnung der Zusicherungen zu den jeweiligen Fehlerinjektionen und Prüfkriterien ist in \cref{tab:testmatrix}
dokumentiert.

\subsection{Test- und Nachweismethoden}
\label{subsec:test_und_nachweismethoden}
Die Auswahl der konkreten Testszenarien orientiert sich primär an der Abdeckung der zuvor definierten Zusicherungen.
Entsprechend wird für jede der neun Zusicherungen systematisch mindestens ein Testszenario definiert.
Die zugrunde liegenden Fehlerklassen werden in \cref{sec:abgrenzungFehlermodell} definiert und durch die Literatur fundiert begründet.
Die vollständige Zuordnung von Systemzusicherung, simulierter Fehlerinjektion und dem dazugehörigen Prüfkriterium ist in
\cref{tab:testmatrix} dokumentiert.
Die Nachweisführung einer Zusicherung kann zusätzlich analytische Anteile umfassen.
Eigenschaften, die sich aus strukturellen Merkmalen der Implementierung ableiten lassen, werden argumentativ anhand dieser
Implementierungsmerkmale begründet.

Komponententests werden im Rahmen der Entwicklung eingesetzt, um einzelne Mechanismen isoliert zu überprüfen und die allgemeine
Funktionsfähigkeit des Systems sicherzustellen.
Diese Tests sind jedoch nicht Bestandteil der Evaluation, deren Fokus auf der Fehlertoleranz liegt und daher
ein unter normalen Bedingungen funktionierendes System voraussetzt.
Stattdessen werden ausschließlich Integrationstests als Evaluationsevidenz herangezogen, die gezielt Fehler injizieren und
dabei das Verhalten und Zusammenspiel der verschiedenen Komponenten überprüfen.
Die konkrete technische Umsetzung der Testumgebung wird in \cref{sec:versuchsaufbau} beschrieben.

\subsection{Fehlerinjektion und Testdurchführung}
\label{subsec:fehlerinjektion_und_testdurchfuehrung}
Die Überprüfung der Zusicherungen erfolgt durch gezielte Herbeiführung der für das jeweilige Testszenario relevanten Fehlerbedingungen.
Ein Testlauf beginnt dabei mit einem definierten Ausgangszustand.
Anschließend wird die zu untersuchende Bedingung erzeugt und das daraus resultierende Systemverhalten beobachtet.
Dieses wird abschließend mit den definierten Prüfkriterien verglichen.

Die Art der Fehlerinjektion wird dabei an das jeweils untersuchte Fehlerszenario angepasst.
Fehler können beispielsweise durch die Unterbrechung von Infrastrukturverbindungen, durch das gezielte Herstellen eines
für Recovery relevanten Systemzustands oder durch das Auslösen definierter Fehler innerhalb einer Verarbeitung erzeugt werden.
Für die Untersuchung von Nebenläufigkeit werden konkurrierende Operationen synchronisiert und gleichzeitig ausgeführt.
Die konkrete technische Umsetzung der jeweiligen Fehlerinjektion und Teststeuerung wird zusammen mit den einzelnen Testszenarien
in \cref{ch:evaluation} beziehungsweise dem Versuchsaufbau in \cref{sec:versuchsaufbau} beschrieben.

\subsection{Reproduzierbarkeit und Auswertung}
\label{subsec:reproduzierbarkeit}
Nebenläufigkeitsabhängige Fehler können ein nicht-deterministisches Verhalten bei Mehrfachausführung desselben Tests verursachen.
Insbesondere Race Conditions hängen vom konkreten Scheduling der beteiligten Threads ab.
Ein einmalig erfolgreicher Testlauf stellt daher für solche Szenarien keinen hinreichenden empirischen Nachweis dar.
Die Tests bezüglich konkurrierender Verarbeitung werden deshalb wiederholt unter identischen Ausgangsbedingungen ausgeführt.
Für jeden Lauf werden dabei dieselben Systemparameter und dieselben Fehlerinjektionsbedingungen verwendet.
Die Anzahl der Wiederholungen wird dabei einheitlich auf 100 festgelegt.
Diese Zahl stellt einen Kompromiss zwischen der Aussagekraft der Untersuchung und dem für die Testdurchführung erforderlichen
Zeitaufwand dar.
Eine höhere Wiederholungszahl erhöht die Wahrscheinlichkeit, selten auftretende Nebenläufigkeitseffekte zu beobachten,
führt jedoch gleichzeitig zu einem höheren Zeitaufwand für die Testdurchführung.
Ausgewertet wird, ob die für das jeweilige Szenario definierten Prüfkriterien über sämtliche Wiederholungen hinweg eingehalten werden.
Bei Tests, bei denen unterschiedliche Ausgänge des Rennens möglich sind, wird zusätzlich überprüft, dass im Verlauf der Wiederholungen
unterschiedliche Ausgänge auftreten.
Dies soll sicherstellen, dass tatsächlich verschiedene Schedules entstehen.
Ein erfolgreicher Durchlauf wird dabei nicht als Beweis für das generelle Ausbleiben eines Fehlers interpretiert.
Die wiederholte Durchführung dient vielmehr dazu, die Robustheit der beobachteten Ergebnisse gegenüber unterschiedlichen
Ausführungsreihenfolgen zu untersuchen.
Dieses Vorgehen bildet zugleich die methodische Grundlage für die Überprüfung von \ref{item:h3}.
Bei den restlichen Tests ist eine Wiederholung nicht notwendig, da diese durchweg deterministisch konstruiert sind und
eine mehrfache Ausführung daher keine unterschiedlichen Schedules abdeckt.
Die konkrete Durchführung, die verwendeten Parameter und die Ergebnisse der Testfälle werden in \cref{ch:evaluation} dokumentiert.
